\DocumentMetadata{
  pdfversion=2.0,
  pdfstandard=ua-2,
  lang=en-US,
  tagging=on,
  tagging-setup={math/setup=mathml-SE,tabsorder=structure}
}
\documentclass[11pt,letterpaper]{article}
\usepackage[margin=0.68in,includefoot]{geometry}
\usepackage{newcomputermodern}
\usepackage{amsmath,amscd,mathtools}
\usepackage{tikz,adjustbox}
\providecommand{\square}{\mdlgwhtsquare}
\usepackage[hidelinks]{hyperref}
\hypersetup{
  pdftitle={Algebra PhD Exam, September 2000},
  pdfauthor={University of Florida Department of Mathematics},
  pdfsubject={Graduate mathematics examination},
  pdfdisplaydoctitle=true,
  bookmarksopen=true
}
\setcounter{secnumdepth}{0}
\setlength{\parindent}{0pt}
\setlength{\parskip}{0.28em}
\setlength{\emergencystretch}{3em}
\setlength{\leftmargini}{2.15em}
\setlength{\leftmarginii}{2.65em}
\setlength{\leftmarginiii}{3em}
\renewcommand{\labelenumi}{\textbf{\theenumi.}}
\pagestyle{empty}
\raggedbottom
\allowdisplaybreaks
\newenvironment{examproblems}[1]{%
  \begin{enumerate}%
  \setcounter{enumi}{#1}%
  \setlength{\topsep}{0.35em}%
  \setlength{\partopsep}{0pt}%
  \setlength{\itemsep}{0.48em}%
  \setlength{\parsep}{0.18em}%
}{\end{enumerate}}
\newenvironment{examsubparts}{%
  \begin{enumerate}%
  \setlength{\topsep}{0.18em}%
  \setlength{\partopsep}{0pt}%
  \setlength{\itemsep}{0.16em}%
  \setlength{\parsep}{0.1em}%
}{\end{enumerate}}
\newenvironment{examinstructions}{%
  \par\begingroup\itshape
}{%
  \par\endgroup\vspace{0.18em}
}
\makeatletter
\renewcommand\section{\@startsection{section}{1}{0pt}%
  {-1.25ex plus -.25ex minus -.1ex}%
  {0.55ex plus .1ex}%
  {\normalfont\large\bfseries}}
\renewcommand{\@maketitle}{%
  \newpage\null\vskip -1em%
  {\footnotesize\bfseries
    \href{https://www.ufl.edu/}{University of Florida}, \href{https://math.ufl.edu/}{Department of Mathematics}\par}%
  \vskip -0.5em%
  \tagstructbegin{tag=Title}%
    {\LARGE\bfseries\href{https://gma.math.ufl.edu/exams/algebra-phd/}{Algebra PhD Exam}, September 2000\par}%
  \tagstructend%
  \vskip 0.25em%
  \tagmcbegin{artifact}%
    \hrule height 1pt%
  \tagmcend%
  \vskip 1em%
}
\makeatother
\title{Algebra PhD Exam, September 2000}
\author{}
\date{}
\begin{document}
\maketitle
\thispagestyle{empty}
\begin{examinstructions}
Work exactly seven out of the following eleven exercises.
\end{examinstructions}
\begin{examproblems}{0}
\item
This problem has three parts.
\begin{examsubparts}
\item[\textnormal{(a)}]
Define solvable group.
\item[\textnormal{(b)}]
Prove that the homomorphic image of a solvable group is solvable.
\item[\textnormal{(c)}]
Prove that a free group is solvable if and only if it is the free group on at most one generator.
\end{examsubparts}
\item
Let~\(G\) be a group; call \(g\in G\) a non-generator if, for each subset~\(X\) of~\(G\) so that \(X\cup\{g\}\) generates~\(G\), then, in fact, \(X\) itself generates~\(G\). Let \(\operatorname{Fr}(G)\) denote the set of all non-generators of~\(G\).
\begin{examsubparts}
\item[\textnormal{(a)}]
Prove that \(\operatorname{Fr}(G)\) is a subgroup of~\(G\).
\item[\textnormal{(b)}]
Show that \(\operatorname{Fr}(G)\) is the intersection of all maximal (proper) subgroups of~\(G\). (Careful with Zorn’s Lemma!)
\end{examsubparts}
\item
Suppose that~\(R\) is a principal ideal domain. Prove that any submodule of a free~\(R\)-module is free.
\item
Suppose \((m,n)=1\). Compute \(\mathbb{Z}_m\otimes_{\mathbb{Z}}\mathbb{Z}_n\). Justify your answer.
\item
Let~\(R\) be a finite semisimple ring with identity. Suppose that no fourth power \(k^4>1\) (\(k\in\mathbb{N}\)) divides \(|R|\). Prove that~\(R\) is commutative, and therefore a direct product of fields.
\item
Suppose that~\(R\) is a ring with identity. Prove that 
\[
\operatorname{Hom}_{\mathrm{Ab}}\left(B,\prod_{i\in I}G_i\right)\cong\prod_{i\in I}\operatorname{Hom}_{\mathrm{Ab}}(B,G_i),
\]
 as right~\(R\)-modules, for all left~\(R\)-modules~\(B\) and all abelian groups \(G_i\) (\(i\in I\)). Ab denotes the category of all abelian groups together with all homomorphisms between them.

\par
You may use resources from category theory; if so, outline your argument so that it is clear which theorems you are appealing to.
\item
Prove Nakayama’s Lemma: let~\(A\) be a commutative ring with identity. Let~\(M\) be a finitely generated~\(A\)-module, and~\(I\) be an ideal of~\(A\), contained in the Jacobson radical \(J(A)\) of~\(A\). Show that if \(IM=M\) then \(M=\{0\}\).
\item
Let~\(A\) be a commutative ring with identity. For each multiplicative system~\(S\) of~\(A\), prove that \(S^{-1}A\) is a flat~\(A\)-module.
\item
Suppose that~\(A\) is a subring of the commutative ring~\(B\) with 1. If~\(B\) is integral over~\(A\), prove that every homomorphism~\(f\) of~\(A\) into the algebraically closed field~\(L\) admits an extension to a homomorphism \(g:B\longrightarrow L\).
\item
Let \(\Theta_p\) be the~\(p\)-th cyclotomic polynomial over the field~\(\mathbb{Q}\), where~\(p\) is a prime number. Show that the Galois group of the splitting field of \(\Theta_p\) over~\(\mathbb{Q}\) is cyclic of order \(p-1\). (Be clear about any theorems you quote.)
\item
Consider the polynomial over the field \(\mathbb{F}_2\) of two elements: \(g(x)=x^4+x^3+x^2+x+1\).
\begin{examsubparts}
\item[\textnormal{(a)}]
Prove that \(g(x)\) is irreducible over \(\mathbb{F}_2\).
\item[\textnormal{(b)}]
Let~\(K\) be a splitting field for \(g(x)\) over \(\mathbb{F}_2\), and let \(r\in K\) be a root of \(g(x)\). Factor \(g(x)\) into irreducibles over \(\mathbb{F}_2(r)\).
\item[\textnormal{(c)}]
Show that \(K=\mathbb{F}_2(r)\).
\item[\textnormal{(d)}]
Find the Galois group \(\operatorname{Gal}(K/\mathbb{F}_2)\).
\end{examsubparts}
\end{examproblems}
\end{document}
